%% kapa: exam=1 round=26 year=2015 period=2 form=A subject=appraisal_law subject_ko=부동산관계법규 seq=15 number=15
%% source: Q-Net 감정평가사 기출문제 (cst00309) artlSeq=5138001 file=extracted/1008/5138001/제26회_감정평가사 1차_기출/제26회_감정평가사 1차_기출/2교시_감평_A형.pdf
%% conversion: kordoc 4.24.51 -> tools/gen1.py (auto); review=proofread-sonnet; flags=none
\begin{question}{15}
국토의 계획 및 이용에 관한 법령상 도시ㆍ군계획시설 부지의 매수 청구에 관한 설명으로 옳은 것은?
\begin{choices}
  \item 매수의무자는 매수 청구를 받은 날부터 1년 이내에 매수 여부를 결정하여 토지소유자에게 알려야 한다.
  \item 지방자치단체인 매수의무자는 부재부동산 소유자의 토지 또는 비업무용 토지로서 매수대금이 2천 만원을 초과하여 그 초과하는 금액을 지급하는 경우에는 도시ㆍ군계획시설채권을 발행하여 지급하여야 한다.
  \item 도시ㆍ군계획시설채권의 구체적인 상환기간은 20년 이내의 범위에서 지방자치단체의 조례로 정한다.
  \item 도시ㆍ군계획시설을 설치하거나 관리하여야 할 의무가 있는 자가 서로 다른 경우에는 관리하여야 할 의무가 있는 자에게 매수 청구하여야 한다.
  \item 매수의무자가 매수하기로 결정한 토지는 매수 결정을 알린 날부터 2년 이내에 매수하여야 한다.
\end{choices}
\end{question}
\answer{5}
